\chapter{State of the Art and Related Work}
\label{ch:related-work}

\Cref{ch:ltn} set out the framework of the thesis and its candidate axioms. This
chapter places them in the literature. It reviews learning-based
detection that carries no explicit knowledge, neurosymbolic work in
cybersecurity and in network intrusion detection, and the benchmarking
practice of the field. It ends with the gap that the thesis addresses.

% ================================================================
% 6.1 Review method
% ================================================================
\section{Review method}
\label{sec:review-method}

The review of neurosymbolic detectors in this chapter is a targeted search. It
is not a systematic review in the sense of PRISMA. The field is young, and the
set of detectors that work on network traffic is small. A protocol with many
screening stages would add cost and little precision. A systematic review
already covers neurosymbolic AI in cybersecurity
\citep{hakim2025nesycybersecurity}, and a survey covers neurosymbolic network
intrusion detection \citep{bizzarri2025survey}. We use them as maps, and we check a paper before we
report its figures.

\subsection{Sources, queries, inclusion and exclusion criteria, time window}
\label{sec:review-sources}

The search ran in September and October 2026 and had two sources. The first is
the literature review of the author's internship, which lists ten papers. The
second is a set of queries to a general web search engine. The terms were
variations of neurosymbolic intrusion detection, logic tensor networks and
network traffic, neuro-symbolic detection with domain knowledge on NSL-KDD,
probabilistic logic programming for intrusion detection, and knowledge graphs
for anomaly detection in cybersecurity. No bibliographic database was queried
directly, and the reference lists of the papers were not followed. A paper that
this search missed is therefore missing from the corpus.

A paper entered the corpus if it met four criteria. It has a neural component
and an explicit symbolic component. It applies them to network traffic or to
data derived from it. It evaluates the method on data. It is published in a
peer-reviewed venue or on arXiv. Work on source code, malware or radio
spectrum was excluded, because the data are not network traffic. Reviews and
position papers were kept as context for \cref{sec:nesy-cybersecurity} and are
not part of the corpus. We set no lower bound on the year, and the earliest paper
of the corpus is from 2022. A paper was also excluded when we could not trace it
to a primary source.

\subsection{Extraction grid}
\label{sec:review-grid}

For every paper of the corpus we recorded seven fields: the approach, the type
of knowledge, the dataset, the baselines, the metrics, the reproducibility
information and the availability of code. The reproducibility field holds the
number of seeds or repetitions and the statistical tests. Three more checks
follow from the questions of the thesis. They ask whether the paper evaluates
real-time operation, whether it reports computational cost and whether it
evaluates its explanations. A field is marked as not stated when the paper does
not say, which is not the same as absent. We also record whether the extraction
rests on the full text or on the abstract alone.

Every statement below that something is not reported holds for the corpus only.
It is not a statement about the field.

% ================================================================
% 6.2 ML and DL for network intrusion detection
% ================================================================
\section{ML and DL for network intrusion detection}
\label{sec:ml-dl-nids}

This section covers detectors that learn from traffic and carry no explicit
knowledge component. They are the counterpart of the neurosymbolic detectors
of the next sections, and the datasets they use are those of the benchmark.
The section relies on surveys and on a few primary papers. The review
protocol is applied to the neurosymbolic detectors and not to this section.

\subsection{Shallow models on NSL-KDD / UNSW-NB15 / CIC-IDS}
\label{sec:shallow-nids}

Surveys of learning-based intrusion detection cover the models that
\cref{sec:classical-models} introduced. \citet{buczak2016survey} give a short
tutorial on each machine learning and data mining method that they review, and
they add recommendations on when to use a given method.
\citet{thakkar2022survey} cover the work of 2008 to 2020. They report that the
methods they survey were tested on datasets such as DARPA, KDD Cup 99 and
NSL-KDD with network features, and they describe CIC-IDS-2017 in detail as a
recent dataset.

The neurosymbolic survey of \citet{bizzarri2025survey} supports the same point
with a table. The table has 19 entries for machine learning or deep learning
approaches that carry no knowledge, and 14 of them are evaluated on KDD99 or
NSL-KDD. The count is ours, and the table is a selection of the authors and not
a census of the field. The authors read it as follows. Deep detectors are
mainly evaluated on synthetic datasets, and on the few recent datasets the
false positive rate is not always reported.

Newer datasets were built to replace the older ones.
\citet{moustafa2015unswnb15} introduce UNSW-NB15 on the ground that KDD98,
KDDCUP99 and NSL-KDD were generated a decade ago and do not fully reflect
current traffic and low footprint attacks. The dataset has nine attack types,
among them Fuzzers, Exploits and Worms.

\subsection[Deep models: CNN, RNN, autoencoder, transformer-based detectors]{Deep models: CNN, RNN, autoencoder,\\ transformer-based detectors}
\label{sec:deep-nids}

\citet{ferrag2020deep} survey deep learning for intrusion detection. They
describe 35 datasets in seven categories and compare recurrent networks, deep
neural networks, restricted Boltzmann machines, deep belief networks,
convolutional networks, deep Boltzmann machines and deep autoencoders. The
comparison runs on CSE-CIC-IDS2018 and Bot-IoT, for binary and multiclass
tasks, and it uses accuracy, false alarm rate and detection rate.

Autoencoders enter the field as anomaly detectors. Kitsune is an ensemble of
autoencoders, called KitNET, that learns online and without supervision
\citep{mirsky2018kitsune}. Each autoencoder reconstructs a subset of the
features of a packet, and an output autoencoder combines the reconstruction
errors into an anomaly score. The model is trained on traffic assumed to be
normal. Its authors contrast this with classification, which they call a
closed-world approach, because a classifier identifies only the classes of its
training set. The term is used here in another sense than the closed world
assumption of \cref{sec:datadriven-not-enough}. Kitsune is evaluated on an IP camera network and on an IoT
network infected with Mirai. The paper reports the true positive rate at fixed
false positive rates, the area under the ROC curve and the equal error rate.
It also reports packet processing rates on a Raspberry Pi. This is the
detector in the case study of \citet{arp2022dosdonts}, where a simple baseline
did better at low false positive rates.

RTIDS is a transformer-based detector \citep{wu2022rtids}. A positional
embedding associates sequential information between features, a stacked
encoder-decoder learns low-dimensional representations, and self-attention
classifies the traffic type. On CICIDS2017 and CIC-DDoS2019 the authors report
F1 scores of 99.17\% and 98.48\%, in a comparison with an SVM and with RNN,
FNN and LSTM models. The features of a flow record have no natural order, as
\cref{sec:dl-architectures} noted for the one-dimensional convolution. A
positional embedding assumes one. We note this as a limit of the design. The remark
does not question the reported scores.

\subsection{The reproducibility and dataset-quality debate}
\label{sec:dataset-quality-debate}

The data behind these results have a history. The KDD Cup 99 data derive from
the 1998 DARPA evaluation \citep{kddcup1999task}, and NSL-KDD is a revision of
them that removes redundant records \citep{cic2009nslkdd}.
\citet{mchugh2000testing} examines the DARPA evaluations of 1998 and 1999. Some
of the methods used are questionable and may have biased the results, and the
evaluators published little on how they validated their test data.
\citet{mahoney2003analysis} compare the simulated background traffic of the
1999 data with real traffic. They find signs of simulation artifacts that could
let an anomaly detector find some attacks with an artificially low false alarm
rate. Their data are the 1999 data and not the 1998
data behind KDD Cup 99. We take the finding as a warning about simulated
traffic and not as a property of NSL-KDD.

\citet{ring2019survey} survey the network-based datasets along 15 properties
grouped in five categories, so that a dataset can be matched to an evaluation
scenario, and they end with recommendations on the use and creation of
datasets. A newer dataset is not a clean one. \citet{engelen2021troubleshooting}
revisit CIC-IDS2017 and find problems in traffic generation, flow construction,
feature extraction and labelling that severely affect it. With an improved
processing methodology they reconstruct or relabel more than 20\% of the
original traffic traces, and benchmarks on the result show significant
improvements. Their study shows how problems in data collection can
change the outcome of an evaluation, and we read this as the reproducibility
side of the debate. The benchmark uses a corrected variant of CIC-IDS2017, and
\cref{ch:implementation} says which one.
% FIXME(open): \cref{ch:implementation} (Chapter 9, Sec. 9.2.2) renders as "??" until that chapter exists; the corrected variant of CIC-IDS2017 is chosen there.

The thesis uses NSL-KDD all the same. We give three reasons and a limit.

The first reason is the content features. R2L and U2R attacks are embedded in
the data portion of a single connection, and the content features of NSL-KDD
are what exposes them to a detector that reads one record
\citep{tavallaee2009kddcup99}. The flow features of CIC-IDS2017 are statistics
computed over each flow \citep{sharafaldin2018cicids2017}, and we note that none
of them is a content feature of the NSL-KDD kind.
Without NSL-KDD, the knowledge about R2L and U2R that \cref{ch:ltn} writes down
could not be tested.

The second reason is comparability. Most of the machine learning and deep
learning entries in the table of Bizzarri et al. use KDD99 or NSL-KDD. The two
detectors built with an \gls{ltn} that we read use KDD99 with the duplicates
removed and CIC-IDS2017 \citep{onchis2022neurosymbolic}, and CIC-IDS2017 alone
\citep{bizzarri2024nsaids}. Experiment E1 runs the pipeline of the thesis on
NSL-KDD under the corrected protocol, on the lineage of data of most published work
(\cref{sec:res-e1}).

The third reason is the role of the dataset. NSL-KDD serves for R2L and U2R and
for E1. The thesis does not draw its general conclusions from
NSL-KDD alone, and it reads a result that rests on content features as specific
to NSL-KDD.

The limit is the one set out above. NSL-KDD is old and synthetic, and a score
that holds only there says nothing about current traffic. The benchmark reports
it as such.

% ================================================================
% 6.3 Neurosymbolic AI in cybersecurity
% ================================================================
\section{Neurosymbolic AI in cybersecurity}
\label{sec:nesy-cybersecurity}

Neurosymbolic work in security is wider than the detection of attacks on
traffic. This section maps it by the way knowledge and learning are joined. The
detectors of network traffic are left to the next section.

\subsection{Surveys and taxonomies}
\label{sec:nesy-cyber-surveys}

\citet{hakim2025nesycybersecurity} review 103 publications on neurosymbolic AI
in cybersecurity, from 2019 to early 2026. They sort the papers in three tiers.
Deep integration covers joint optimisation, as in logic tensor networks.
Structured interaction covers systems whose neural and symbolic parts exchange
results, for example a knowledge graph that enriches a neural encoder or a
language model paired with a solver. The third tier holds purely statistical
systems kept for comparison. The authors name the lack of standard frameworks
for evaluating these systems as the most critical gap of the field, and they
note that the datasets in use carry no symbolic knowledge. They also warn that
the satisfiability computation of systems built on logic tensor networks may
add an inference latency that does not fit wire speed. In our reading the warning holds for a system that evaluates the formulas at
inference. In the design of the thesis the
formulas act in the loss only, as \cref{sec:inference-vs-training} explains, and
the cost at inference is a question for the experiments and not an assumption.

\citet{bizzarri2025survey} survey neurosymbolic AI for network intrusion
detection. They find that latency and throughput are not consistently reported
and that the overhead of joining symbolic and neural processing is rarely
quantified. Two shorter papers state a programme. \citet{grov2024neurosymbolic}
argue for neurosymbolic cyber defence and give use cases with proof-of-concept
experiments. \citet{piplai2023knowledge} present knowledge graphs as the
symbolic side of neurosymbolic AI for security and privacy, with explainability
and safety as the aim. We read both as position papers, and we use them for
framing and not as evidence.

\subsection{Knowledge-graph based approaches}
\label{sec:nesy-kg}

A knowledge graph holds the context of a detection, such as assets, services and
the relations between them. \citet{tailhardat2025anomaly} survey anomaly
detection with knowledge graphs for network management and cybersecurity. Their
text is a preprint at the time of writing.
\citet{garrido2021knowledge} apply relational learning to a knowledge graph of
an industrial automation setting, where network traffic analysed with Zeek is
one of the data sources, to put alerts in context. The model is a probabilistic
tensor factorisation and not a neural network, so the work sits at the edge of
the field.

In both, knowledge enters beside or after the detector and not in its training.
This is closest to the post-hoc family of \cref{sec:post-hoc-verification}, where
the model is left as it is. A detector of this kind appears in the next section.

\subsection{Logic-as-loss approaches, LTN applications}
\label{sec:nesy-logic-loss}

The thesis belongs to the family that puts formulas in the loss. In the corpus
of this chapter, two detectors are of this family. Both are built with logic
tensor networks and both work on network traffic, so they are read in
\cref{sec:nesy-nids}. No detector of the corpus uses a semantic loss or DL2. The
general literature on constraints in the loss does not concern security in
particular and is covered in \cref{sec:logic-as-loss}.

\subsection{Probabilistic logic programming approaches}
\label{sec:nesy-plp}

The preprint of Bizzarri et al. presents DeepProbLog, DeepStochLog and NeurASP
as methods, and it reports no application of them to intrusion detection on
benchmark data \citep{bizzarri2024synergistic}. We found no detector built with
them either.

The probabilistic logic of the corpus is Markov logic. KnowGraph writes domain
knowledge as weighted first-order formulas of three kinds, attribute rules,
hierarchy rules and exclusion rules, and reasons over them on graph data with
variational inference \citep{zhou2024knowgraph}. The data are authentication
events of the LANL set, 1.6 billion over 58 days, and a marketplace set. It is
not a classifier of traffic. In our view its three kinds of rule correspond to
the groups of axioms of \cref{sec:ddos-axioms}, which gives the thesis a point
of comparison in another logic. KnowGraph reports the area under the ROC curve, the average
precision and the true positive rate at fixed false positive rates. We found no
latency or cost figures in it.

\subsection{LLM and formal verification hybrids}
\label{sec:nesy-llm-formal}

Hybrids of language models and formal verification have been reviewed for
software vulnerabilities \citep{tihanyi2025vulnerability}. In our simplified
description, a model proposes and a verifier checks. The artefacts are code and specifications, which are discrete
and well formed, and a flow record is neither. Hakim et al. place such pairings
in the structured interaction tier. We found no use of the scheme on network
traffic. It is outside the scope of the thesis, whose detectors contain no
language model.

% ================================================================
% 6.4 Neurosymbolic AI applied specifically to NIDS
% ================================================================
\section{Neurosymbolic AI applied specifically to NIDS}
\label{sec:nesy-nids}

\subsection{Findings of the internship literature review}
\label{sec:internship-findings}

The internship of the author produced a review of ten papers. We used it as a
map and checked every paper against its primary source. Three papers are
detectors that we could confirm: \citet{onchis2022neurosymbolic},
\citet{bizzarri2024nsaids} and \citet{kalutharage2025neurosymbolic}. Three are
reviews, of which \citet{hakim2025nesycybersecurity} is the one on
cybersecurity. One concerns radio spectrum management, which is not network
traffic classification. For three we could not find a primary source. The
bibliography of the review also gave wrong venues for the papers that we did
confirm. We corrected the venues and dropped the papers that we could not trace.

The search of \cref{sec:review-sources} added five works. The corpus then has
eight detectors, listed in \cref{tab:nesy-nids-grid}. They are
\citet{tran2025odxu}, \citet{tran2025transfer}, \citet{aresrobledo2025hyperparameter},
\citet{zhou2024knowgraph} and \citet{alemayehu2026deployment}, in addition to the
three above.

\begin{table}[htbp]
\centering
\footnotesize
\setlength{\tabcolsep}{3pt}
\caption{The corpus of neurosymbolic detectors. The first column gives the
reference and the other four are the fields of the extraction grid. An
asterisk marks an extraction based on the abstract alone.}
\label{tab:nesy-nids-grid}
\begin{tabular}{@{}p{2.3cm}p{3.2cm}p{2.6cm}p{1.9cm}p{2.0cm}@{}}
\toprule
Work & Symbolic part & Data & Baselines & Metrics \\
\midrule
\citet{onchis2022neurosymbolic}
  & Class and exclusivity axioms in an LTN
  & KDD99 deduplicated, one afternoon of CIC-IDS2017
  & DNN
  & Accuracy, KB satisfaction \\
\citet{bizzarri2024nsaids}
  & Benign against attack axiom in an LTN, dropped at test
  & CIC-IDS2017 packets
  & 1D-CNN, two LTN variants
  & Accuracy, F1 \\
\citet{tran2025odxu}
  & Decision-path rules of XGBoost on learned embeddings
  & CIC-IDS2017, ACI-IoT-2023
  & DNN, 1D-CNN
  & Accuracy, false omission rate, AUROC, training time \\
\citet{tran2025transfer}
  & Same, with transfer learning
  & ACI-IoT-2023, CIC-IDS2017
  & DNN, 1D-CNN
  & Same \\
\citet{kalutharage2025neurosymbolic}
  & Knowledge graph that validates alerts
  & USB-IDS, own IoT capture
  & DT, RF, DNN from prior work
  & Accuracy \\
\citet{zhou2024knowgraph}
  & Weighted first-order rules, Markov logic
  & LANL, marketplace
  & Euler, EncG, GCN, DGI
  & AUC, AP, TPR at fixed FPR \\
\citet{aresrobledo2025hyperparameter}$^{*}$
  & Symbolic reasoning in a graph pipeline
  & IoTID20, UNSW-NB15
  & not stated
  & MCC \\
\citet{alemayehu2026deployment}
  & Shallow decision tree with a GRU
  & CIC-DDoS2019, Edge-IIoTset, CICIoT23
  & GRU and tree alone, five classical models
  & Accuracy, MCC, FPR, FNR, latency, memory \\
\bottomrule
\end{tabular}
\end{table}

The label neurosymbolic covers different things in these papers. In the two
\gls{ltn} detectors, formulas enter the loss. In KnowGraph, weighted rules form
a reasoning layer over a graph. In the work of Kalutharage et al. a knowledge
graph checks the alerts of an autoencoder after the detection. In the two papers
of Tran et al. and in the paper of Alemayehu et al. the symbolic part is a tree
model learned from data. Ares-Robledo et al. add symbolic reasoning to an
unsupervised graph pipeline. Only in the first group is the network that
classifies the traffic trained directly to satisfy the formulas. KnowGraph and
the work of Kalutharage et al. also use knowledge written by people. In
KnowGraph the detection models are trained without the rules. The weighted rules
then define a Markov logic network, and a separate reasoning network is trained
to approximate it. The rules therefore act on training through this second
network, which is kept at prediction time \citep{zhou2024knowgraph}. In the work
of Kalutharage et al. the
experts set thresholds in a graph that acts only after the detection.
This classification is ours. It means that the benchmark of the thesis matches the
first group, and it does not match the others.

\subsection{Critical reading: what is actually measured, what is not}
\label{sec:critical-reading}

We start with the two detectors that put formulas in the loss.

\citet{onchis2022neurosymbolic} train on KDD99 after the removal of duplicates,
which cuts it from 494{,}020 to 145{,}585 records, and on one afternoon of
CIC-IDS2017. On KDD99 the paper has two classifiers. The first predicts the
protocol type and the connection status and serves to show queries. The
knowledge base of the intrusion classifier states that each class holds for its
examples and that the classes exclude each other. The CIC-IDS2017 classifier has
the three class axioms only. No axiom of the two intrusion classifiers mentions a
property of the traffic. The axioms constrain labels. The satisfaction of such a
base follows the quality of the classifier, as the authors say themselves when
they note that a wrong prediction lowers the truth values. It is reported on
training and test data. In our reading it differs from accuracy because it also
depends on the confidence of the outputs. On KDD99 the accuracy reaches 0.995 on training and 0.995 on test data
against 0.989 and 0.987 for a dense network, while the satisfaction on test data
stays near 0.86, as we read the figure. On CIC-IDS2017 the authors report an
accuracy of 0.95 and a satisfaction of about 0.85.

The comparison with the dense network is not matched. The authors keep the
number of hidden layers and of neurons, but the LTN uses a sigmoid output and
elu units and the baseline a softmax and relu units. The paper reports no
repetitions, seeds or statistical tests. On
CIC-IDS2017 the figures, as we read them, show a validation accuracy of the
dense network above that of the LTN late in training, with larger swings. The
argument of the paper is about the stability of the curve and not about the
level. No precision, recall or per-class result is reported, although U2R has a
few dozen records in the data, as a figure of the paper shows. The conclusions say that the symbolic part speeds
up training and reduces overfitting. The paper gives no training times, and the
overfitting argument rests on one training curve per model. The authors also present
CIC-IDS2017 as a dataset that aims to solve the faults of the earlier ones, while
\citet{engelen2021troubleshooting} find faults of its own. The article links the datasets and
not the code.

\citet{bizzarri2024nsaids} work on CIC-IDS2017 packets, with the first 1500
payload bytes as features and a 1D-CNN as the network. The knowledge base has one
pair of axioms, benign examples are benign and attacks are not. Attack classes
are held out of training to stand for unknown attacks. The text speaks of five
such classes, while the table and the results of the paper list six. After 50
epochs, on the nine known classes the LTN reaches 81.08\% accuracy against
80.99\% for the same CNN. On the
unknown attacks, seen as a binary problem, it reaches 60.47\% against 48.34\%,
and an F1 score of 75.37\% against 65.18\%. Two points matter for this thesis.
The baseline is the same network trained without the logic term, which is the
matched capacity that \cref{sec:bias-variance-generalization} requires. And the model
used at test time is the bare CNN, as in our own design. The extraction found
no seeds, no test of significance, no cost figure and no satisfaction
measurement. The absolute accuracy on unknown attacks is low, and the gain
rests on one split.

The other detectors measure other things. Both papers of
\citet{tran2025odxu,tran2025transfer} report accuracy, false omission rate and
uncertainty measures against a dense network and a 1D-CNN. Both give
training times and the hardware. The two papers report the same transfer
experiments, with the same baseline scores and training times, so we do not
count them as independent evidence. Neither gives seeds or tests, and the
symbolic part carries no prior knowledge. \citet{kalutharage2025neurosymbolic}
compare a knowledge graph built from data with one enhanced by experts, and the
overall accuracy is 0.91 against 0.97. This is the only comparison of the
quality of the knowledge in the corpus. The mapping to MITRE ATT\&CK is said to
be 100\% accurate and was confirmed by manual validation, with no count of
instances. The detector runs on a Raspberry Pi in the description, and no latency
or memory is given. KnowGraph reports the true positive rate at fixed false
positive rates, which ties its results to the number of false alarms. By their abstract,
\citet{aresrobledo2025hyperparameter} report an MCC of 0.97 on IoTID20 and 0.91 on
UNSW-NB15. \citet{alemayehu2026deployment} measure the cost in most detail. Model
inference takes 0.58 to 0.79 milliseconds per sample on a workstation CPU, and
the authors also report latency, memory and energy for a deployment on an
industrial gateway and on a Raspberry Pi. They retrain the neural part over ten seeds and give
confidence intervals. On Edge-IIoTset the decision tree alone reaches 100\%
accuracy, and the authors themselves call that dataset linearly separable and
treat it as a sanity check.

Across the corpus, four of the eight detectors use CIC-IDS2017, one uses KDD99
and none uses NSL-KDD. Accuracy is the headline metric of the two detectors with
logic in the loss. Seven extractions rest on the full text. Of these, only the
paper of Alemayehu et al. repeats runs over seeds, and only it mentions a
released repository, without giving a link. None tests the difference between
two detectors, and none reports the violation of rules by the predictions. Only
Onchis et al. report the satisfaction of a knowledge base on test data, and their
axioms constrain the labels only. Cost
appears as training time in the two papers of Tran et al. and as latency,
memory and energy in the paper of Alemayehu et al. Explanations are evaluated by
manual validation in one paper, and qualitatively or not at all in the rest. These statements hold for
the eight papers. The corpus supports the claim that logic in the loss can work,
and it does not tell when it pays.

% ================================================================
% 6.5 Benchmarking methodology for learning-based security systems
% ================================================================
\section{Benchmarking methodology for learning-based security systems}
\label{sec:benchmarking-practice}

The studies above differ in protocol, in metric and in rigour. This section
collects the practice that a fair comparison can draw on. It comes from
security, from the evaluation of classifiers and from the benchmarking of
machine learning systems. The thesis fixes its own choices in
\cref{ch:benchmarking-methodology}.

\subsection{Evaluation protocols: temporal splits, cross-attack, leave-one-attack-out}
\label{sec:evaluation-protocols}

\citet{pendlebury2019tesseract} show how a split inflates results. They name
two flaws in malware classification. Spatial bias arises when training and test
data do not reflect the distribution met in deployment. Temporal bias arises
when the split ignores time and produces impossible configurations. They
evaluate under space and time constraints, and on 129{,}000 Android
applications spanning three years they find that three published classifiers had
been evaluated with bias. The data are malware and not traffic, but the
argument that a random split mixes earlier and later samples carries over.

For traffic, \citet{layeghy2022generalisability} train models on one of four
benchmark datasets and test them on the other three. Scores are near perfect on
the dataset of training and fall sharply elsewhere. No model generalises over
all four, and the result changes when source and target are swapped.
In F1 score, unsupervised models lose 28.33\% on average and supervised ones
56.28\%.
\citet{sarhan2021zeroshot} leave out one attack class at a time to simulate a
zero-day attack, and they measure how often the held-out class is detected.
The detector of \citet{bizzarri2024nsaids} holds out several classes at once. The
stress scenarios of the thesis for unseen attacks and for transfer between
datasets follow these protocols.
% FIXME(open): the split strategy, temporal and random variants, is defined in Ch. 9, Sec. 9.2.4; check that it matches this section.

\subsection{Metrics beyond accuracy: MCC, PR-AUC, FPR@TPR, alert budget, cost curves}
\label{sec:metrics-beyond-accuracy}

\citet{chicco2020mcc} argue that the Matthews correlation coefficient is high
only if the prediction is good in all four categories of the confusion matrix,
while accuracy and the F1 score can be inflated on imbalanced data.
\citet{saito2015precision} show that the precision-recall curve changes with the
class ratio and the ROC curve does not, which makes the first more informative
when the classes are imbalanced.

Security papers often report the true positive rate at a fixed false positive
rate. Kitsune does so \citep{mirsky2018kitsune}, and so does KnowGraph
\citep{zhou2024knowgraph}. A fixed rate speaks to the base-rate problem of
\cref{sec:base-rate-false-positives}, because it ties the comparison to the
number of false alarms. We call the alert budget
the number of alerts per period that analysts can handle. With the volume of
traffic it fixes the largest false positive rate that a deployment accepts, and
the true positive rate at that rate is the figure to compare. The definition
is ours. \citet{drummond2006cost} give a third view. Cost curves show the
expected cost of a classifier over class distributions and misclassification
costs, and the authors argue that they suit several uses better than ROC
curves, among them confidence intervals and the comparison of classifiers.
Prevalence in traffic is not fixed, and a cost curve shows how a ranking of
detectors changes with it. We propose this use for the benchmark.

\subsection{Systems-benchmarking practice (MLPerf-style methodology) transplanted here}
\label{sec:mlperf-practice}

\citet{reddi2020mlperf} present MLPerf Inference as a set of rules and best
practices that make results comparable across systems of very different
architecture. Its first round gave more than 600 reproducible measurements from
14 organisations. \citet{mattson2020mlperf} measure training as the time to
reach a target quality and not as throughput. They give two reasons. An
optimisation that raises throughput can raise the time to solution, and
training is stochastic, so the time to solution varies a lot. They add that
the diversity of systems makes a comparison with one binary and fixed
hyperparameters difficult.

Three habits carry over to a detector benchmark. The quality target is fixed
before the measurement. Runs are repeated and the spread is reported. The rules
on what may be tuned are written down. The suites themselves do not carry over.
They target deep learning workloads on accelerators, and neither covers a
tabular detector or a logic term. The thesis takes the discipline and not the
suites.

\subsection{Statistical comparison of classifiers across datasets}
\label{sec:statistical-comparison}

\citet{demsar2006statistical} recommends the Wilcoxon signed-ranks test for two
classifiers on several datasets, and the Friedman test with post-hoc tests for
more, with critical difference diagrams to show the result. The benchmark has
three datasets. With three datasets the smallest two-sided p-value of the
signed-ranks test is 0.25, so no difference can be called significant at the
5\% level. This follows from the exact distribution of the test, and we did the
calculation ourselves. A comparison across the datasets can then be descriptive
only, and the tests belong to repeated splits within each dataset.
\citet{bouthillier2021variance} name three sources of variance, data sampling,
parameter initialisation and hyperparameter choice. They show that an estimator
that varies several of them at once comes close to the ideal one at a lower
cost. Repeated runs on one dataset are not independent samples of datasets, so
a test on them supports a statement about that dataset. In the corpus of the
previous section, one paper repeats its runs over seeds and none tests a
difference between detectors.

% ================================================================
% 6.6 Synthesis table
% ================================================================
\section{Synthesis of the related work}
\label{sec:synthesis}

\Cref{tab:synthesis} puts the corpus on the axes of the thesis. A dash marks a
field that the extraction did not find, and n.s. marks a field that the source
does not state. An asterisk marks an extraction based on the abstract.

\begin{table}[htbp]
\centering
\footnotesize
\setlength{\tabcolsep}{3pt}
\caption{Related work against the questions of the thesis. The classification
of the family is ours.}
\label{tab:synthesis}
\begin{tabular}{@{}p{2.2cm}p{2.2cm}p{2.3cm}p{1.5cm}p{1.5cm}p{1.9cm}@{}}
\toprule
Work & Family & Data & Real time & Cost & Explanation \\
\midrule
\citet{onchis2022neurosymbolic}
  & Logic in the loss, LTN & KDD99, CIC-IDS2017 & --  & -- & Queries, qualitative \\
\citet{bizzarri2024nsaids}
  & Logic in the loss, LTN & CIC-IDS2017 & -- & -- & -- \\
\citet{tran2025odxu}
  & Rules from trees & CIC-IDS2017, ACI-IoT-2023 & -- & Training time, hardware & -- \\
\citet{tran2025transfer}
  & Rules from trees & ACI-IoT-2023, CIC-IDS2017 & -- & Training time, hardware & -- \\
\citet{kalutharage2025neurosymbolic}
  & Graph after detection & USB-IDS, IoT capture & Claimed, not measured & -- & Manual validation \\
\citet{zhou2024knowgraph}
  & Markov logic on graphs & LANL, marketplace & -- & -- & -- \\
\citet{aresrobledo2025hyperparameter}$^{*}$
  & Symbolic graph pipeline & IoTID20, UNSW-NB15 & n.s. & n.s. & n.s. \\
\citet{alemayehu2026deployment}
  & Tree with a GRU & CIC-DDoS2019 and two more & Latency on edge devices & Latency, memory, energy & -- \\
\bottomrule
\end{tabular}
\end{table}

Only one source mentions a released repository, and it gives no link, so the
column on code is left out. Three papers measure a cost, two as training time
and one as latency, memory and energy. The paper that reports latency is the
only one that uses CIC-DDoS2019. The two detectors with logic in the loss report neither cost nor
real-time behaviour. The table does not say that the missing figures were not
measured, only that the papers do not show them.

% ================================================================
% 6.7 Research gap
% ================================================================
\section{Research gap addressed by this thesis}
\label{sec:research-gap}

The literature shows that logic can be put in the loss of a traffic classifier,
and in the best case it gives a gain on attacks never seen in training. It
does not show when that pays, and six gaps follow from the reading above.

The first concerns the knowledge. In the two detectors with logic in the loss,
the axioms constrain labels and say nothing about the traffic. The thesis
encodes conditions on flows, windows and sources, and tests whether they help.
The second concerns the comparison. Among the detectors with logic in the
loss, only one trains the same network without the logic term. The thesis
matches architecture, features and tuning budget for every pair. The third
concerns the data. The two detectors with logic in the loss use one or two
datasets each, and neither tests a fixed model on a dataset it was not trained
on. The thesis uses three datasets, with stress scenarios for scarce
labels, imbalance, unseen attacks, transfer and degraded knowledge.

The fourth is reproducibility. One source of the corpus repeats runs over seeds,
none tests the difference between two detectors, and none links its code. The
survey of Hakim et al. calls for standard evaluation frameworks. The thesis
publishes its pipeline and reports repeated runs with paired tests. The fifth is
cost. Training time appears in two papers, and latency, memory and energy in
one, which has no logic in the loss.
The thesis measures training and inference cost in batch and in streaming
operation. The sixth is usability. No paper reports the effort of writing the
knowledge or evaluates the explanations beyond a manual check. The thesis
measures both. On the knowledge itself, one paper reports the satisfaction of a
knowledge base, with axioms on labels, and none reports how often the predictions
violate a rule. The two measures of C5 extend that practice to rules on traffic
features, one rule at a time and for every arm.

These gaps map to the research questions. The first two concern RQ1 and RQ2,
the third and fourth RQ5, the fifth RQ3 and the sixth RQ4. Every statement
holds for the corpus of \cref{sec:review-method}, and a paper that the search
missed could narrow a gap.

% ================================================================
% 6.8 Summary
% ================================================================
\section{Summary}
\label{sec:ch6-summary}

Detectors without knowledge are evaluated mostly on old and synthetic data or on
data with known flaws, and the literature on datasets explains why the results
cannot be taken at face value. Neurosymbolic work in cybersecurity is varied. In
network traffic it comes to eight detectors, of which two put formulas in the
loss. They report accuracy and rarely report cost. Only one repeats its runs
over seeds, and none tests a difference between detectors or links its code. The practice of benchmarking offers protocols, metrics and
statistical care that the thesis adopts. What the literature leaves open is
when logic in the loss pays, under a matched comparison, on several datasets and
with its cost measured. The next chapter turns from the
literature to the use cases and the design of the system.
